\documentclass[11pt,twoside,final,openright]{xenstyle}
\usepackage{a4,graphicx,setspace,times}
\usepackage{comment,parskip}
\setstretch{1.15}

\begin{document}

% TITLE PAGE
\pagestyle{empty}
\begin{center}
\vspace*{\fill}
\includegraphics{figs/xenlogo.eps}
\vfill
\vfill
\vfill
\begin{tabular}{l}
{\Huge \bf Interface manual} \\[4mm]
{\huge Xen v2.0 for x86} \\[80mm]

{\Large Xen is Copyright (c) 2002-2004, The Xen Team} \\[3mm]
{\Large University of Cambridge, UK} \\[20mm]
\end{tabular}
\end{center}

{\bf
DISCLAIMER: This documentation is currently under active development
and as such there may be mistakes and omissions --- watch out for
these and please report any you find to the developer's mailing list.
Contributions of material, suggestions and corrections are welcome.
}

\vfill
\cleardoublepage

% TABLE OF CONTENTS
\pagestyle{plain}
\pagenumbering{roman}
{ \parskip 0pt plus 1pt
  \tableofcontents }
\cleardoublepage

% PREPARE FOR MAIN TEXT
\pagenumbering{arabic}
\raggedbottom
\widowpenalty=10000
\clubpenalty=10000
\parindent=0pt
\parskip=5pt
\renewcommand{\topfraction}{.8}
\renewcommand{\bottomfraction}{.8}
\renewcommand{\textfraction}{.2}
\renewcommand{\floatpagefraction}{.8}
\setstretch{1.15}

\chapter{Introduction}

Xen allows the hardware resources of a machine to be virtualized and
dynamically partitioned, allowing multiple different {\em guest}
operating system images to be run simultaneously.  Virtualizing the
machine in this manner provides considerable flexibility, for example
allowing different users to choose their preferred operating system
(e.g., Linux, NetBSD, or a custom operating system).  Furthermore, Xen
provides secure partitioning between virtual machines (known as
{\em domains} in Xen terminology), and enables better resource
accounting and QoS isolation than can be achieved with a conventional
operating system. 

Xen essentially takes a `whole machine' virtualization approach as
pioneered by IBM VM/370.  However, unlike VM/370 or more recent
efforts such as VMWare and Virtual PC, Xen does not attempt to
completely virtualize the underlying hardware.  Instead parts of the
hosted guest operating systems are modified to work with the VMM; the
operating system is effectively ported to a new target architecture,
typically requiring changes in just the machine-dependent code.  The
user-level API is unchanged, and so existing binaries and operating
system distributions work without modification.

In addition to exporting virtualized instances of CPU, memory, network
and block devices, Xen exposes a control interface to manage how these
resources are shared between the running domains. Access to the
control interface is restricted: it may only be used by one
specially-privileged VM, known as {\em domain 0}.  This domain is a
required part of any Xen-based server and runs the application software
that manages the control-plane aspects of the platform.  Running the
control software in {\it domain 0}, distinct from the hypervisor
itself, allows the Xen framework to separate the notions of 
mechanism and policy within the system.



\chapter{Virtual Architecture}

On a Xen-based system, the hypervisor itself runs in {\it ring 0}.  It
has full access to the physical memory available in the system and is
responsible for allocating portions of it to the domains.  Guest
operating systems run in and use {\it rings 1}, {\it 2} and {\it 3} as
they see fit. Segmentation is used to prevent the guest OS from
accessing the portion of the address space that is reserved for
Xen. We expect most guest operating systems will use ring 1 for their
own operation and place applications in ring 3.

In this chapter we consider the basic virtual architecture provided 
by Xen: the basic CPU state, exception and interrupt handling, and
time. Other aspects such as memory and device access are discussed 
in later chapters. 

\section{CPU state}

All privileged state must be handled by Xen.  The guest OS has no
direct access to CR3 and is not permitted to update privileged bits in
EFLAGS. Guest OSes use \emph{hypercalls} to invoke operations in Xen; 
these are analagous to system calls but occur from ring 1 to ring 0. 

A list of all hypercalls is given in Appendix~\ref{a:hypercalls}. 



\section{Exceptions}

A virtual IDT is provided --- a domain can submit a table of trap
handlers to Xen via the {\tt set\_trap\_table()} hypercall.  Most trap
handlers are identical to native x86 handlers, although the page-fault
handler is somewhat different.


\section{Interrupts and events}

Interrupts are virtualized by mapping them to events, which are
delivered asynchronously to the target domain.  A guest OS can map
these events onto its standard interrupt dispatch mechanisms.  Xen 
is responsible for determining the target domain that will handle 
each physical interrupt source. 



\section{Time}

Guest operating systems need to be aware of the passage of both real
(or wallclock) time and their own `virtual time' (the time for
which they have been executing). Furthermore, Xen has a notion of 
time which is used for scheduling. The following notions of 
time are provided: 

\begin{description}
\item[Cycle counter time.]

This provides a fine-grained time reference.  The cycle counter time is
used to accurately extrapolate the other time references.  On SMP machines
it is currently assumed that the cycle counter time is synchronised between
CPUs.  The current x86-based implementation achieves this within inter-CPU
communication latencies.

\item[System time.]

This is a 64-bit counter which holds the number of nanoseconds that
have elapsed since system boot.


\item[Wall clock time.]

This is the time of day in a Unix-style {\tt struct timeval} (seconds
and microseconds since 1 January 1970, adjusted by leap seconds).  An
NTP client hosted by {\it domain 0} can keep this value accurate.  


\item[Domain virtual time.]

This progresses at the same pace as system time, but only while a
domain is executing --- it stops while a domain is de-scheduled.
Therefore the share of the CPU that a domain receives is indicated by
the rate at which its virtual time increases.

\end{description}


Xen exports timestamps for system time and wall-clock time to guest
operating systems through a shared page of memory.  Xen also provides
the cycle counter time at the instant the timestamps were calculated,
and the CPU frequency in Hertz.  This allows the guest to extrapolate
system and wall-clock times accurately based on the current cycle
counter time.

Since all time stamps need to be updated and read \emph{atomically}
two version numbers are also stored in the shared info page. The 
first is incremented prior to an update, while the second is only
incremented afterwards. Thus a guest can be sure that it read a consistent 
state by checking the two version numbers are equal. 

Xen includes a periodic ticker which sends a timer event to the
currently executing domain every 10ms.  The Xen scheduler also sends a
timer event whenever a domain is scheduled; this allows the guest OS
to adjust for the time that has passed while it has been inactive.  In
addition, Xen allows each domain to request that they receive a timer
event sent at a specified system time by using the {\tt
set\_timer\_op()} hypercall.  Guest OSes may use this timer to
implement timeout values when they block.


\chapter{Memory}
\label{c:memory} 

Xen is responsible for managing the allocation of physical memory to
domains, and for ensuring safe use of the paging and segmentation
hardware.


\section{Memory Allocation}


Xen resides within a small fixed portion of physical memory; it also
reserves the top 64MB of every virtual address space. The remaining
physical memory is available for allocation to domains at a page
granularity.  Xen tracks the ownership and use of each page, which
allows it to enforce secure partitioning between domains.

Each domain has a maximum and current physical memory allocation. 
A guest OS may run a `balloon driver' to dynamically adjust its 
current memory allocation up to its limit. 


%% XXX SMH: I use machine and physical in the next section (which 
%% is kinda required for consistency with code); wonder if this 
%% section should use same terms? 
%%
%% Probably. 
%%
%% Merging this and below section at some point prob makes sense. 

\section{Pseudo-Physical Memory}

Since physical memory is allocated and freed on a page granularity,
there is no gaurantee that a domain will receive a contiguous stretch
of physical memory. However most operating systems do not have good
support for operating in a fragmented physical address space. To aid
porting such operating systems to run on top of Xen, we make a
distinction between \emph{machine memory} and \emph{pseduo-physical
memory}.

Put simply, machine memory refers to the entire amount of memory
installed in the machine, including that reserved by Xen, in use by
various domains, or currently unallocated. We consider machine memory
to comprise a set of 4K \emph{machine page frames} numbered
consecutively starting from 0. Machine frame numbers mean the same
within Xen or any domain.

Pseudo-physical memory, on the other hand, is a per-domain
abstraction. It allows a guest operating system to consider its memory
allocation to consist of a contiguous range of physical page frames
starting at physical frame 0, despite the fact that the underlying
machine page frames may be sparsely allocated and in any order.

To achieve this, Xen maintains a globally readable {\it
machine-to-physical} table which records the mapping from machine page
frames to pseudo-physical ones. In addition, each domain is supplied
with a {\it physical-to-machine} table which performs the inverse
mapping. Clearly the machine-to-physical table has size proportional
to the amount of RAM installed in the machine, while each
physical-to-machine table has size proportional to the memory
allocation of the given domain.

Architecture dependent code in guest operating systems can then use
the two tables to provide the abstraction of pseudo-physical
memory. In general, only certain specialized parts of the operating
system (such as page table management) needs to understand the
difference between machine and pseudo-physical addresses.

\section{Page Table Updates}

In the default mode of operation, Xen enforces read-only access to
page tables and requires guest operating systems to explicitly request
any modifications.  Xen validates all such requests and only applies
updates that it deems safe.  This is necessary to prevent domains from
adding arbitrary mappings to their page tables.

To aid validation, Xen associates a type and reference count with each
memory page. A page has one of the following
mutually-exclusive types at any point in time: page directory ({\sf
PD}), page table ({\sf PT}), local descriptor table ({\sf LDT}),
global descriptor table ({\sf GDT}), or writable ({\sf RW}). Note that
a guest OS may always create readable mappings of its own memory 
regardless of its current type. 
%%% XXX: possibly explain more about ref count 'lifecyle' here?
This mechanism is used to
maintain the invariants required for safety; for example, a domain
cannot have a writable mapping to any part of a page table as this
would require the page concerned to simultaneously be of types {\sf
  PT} and {\sf RW}.


%\section{Writable Page Tables}

Xen also provides an alternative mode of operation in which guests be
have the illusion that their page tables are directly writable.  Of
course this is not really the case, since Xen must still validate
modifications to ensure secure partitioning. To this end, Xen traps
any write attempt to a memory page of type {\sf PT} (i.e., that is
currently part of a page table).  If such an access occurs, Xen
temporarily allows write access to that page while at the same time
{\em disconnecting} it from the page table that is currently in
use. This allows the guest to safely make updates to the page because
the newly-updated entries cannot be used by the MMU until Xen
revalidates and reconnects the page.
Reconnection occurs automatically in a number of situations: for
example, when the guest modifies a different page-table page, when the
domain is preempted, or whenever the guest uses Xen's explicit
page-table update interfaces.


\section{Segment Descriptor Tables}

On boot a guest is supplied with a default GDT, which does not reside
within its own memory allocation.  If the guest wishes to use other
than the default `flat' ring-1 and ring-3 segments that this GDT
provides, it must register a custom GDT and/or LDT with Xen,
allocated from its own memory. Note that a number of GDT 
entries are reserved by Xen -- any custom GDT must also include
sufficent space for these entries. 

For example, the following hypercall is used to specify a new GDT: 

\begin{quote}
int {\bf set\_gdt}(unsigned long *{\em frame\_list}, int {\em entries})

{\em frame\_list}: An array of up to 16 machine page frames within
which the GDT resides.  Any frame registered as a GDT frame may only
be mapped read-only within the guest's address space (e.g., no
writable mappings, no use as a page-table page, and so on).

{\em entries}: The number of descriptor-entry slots in the GDT.  Note
that the table must be large enough to contain Xen's reserved entries;
thus we must have `{\em entries $>$ LAST\_RESERVED\_GDT\_ENTRY}\ '.
Note also that, after registering the GDT, slots {\em FIRST\_} through
{\em LAST\_RESERVED\_GDT\_ENTRY} are no longer usable by the guest and
may be overwritten by Xen.
\end{quote}

The LDT is updated via the generic MMU update mechanism (i.e., via 
the {\tt mmu\_update()} hypercall. 

\section{Start of Day} 

The start-of-day environment for guest operating systems is rather
different to that provided by the underlying hardware. In particular,
the processor is already executing in protected mode with paging
enabled.

{\it Domain 0} is created and booted by Xen itself. For all subsequent
donains, the analogue of the boot-loader is the {\it domain builder},
user-space software running in {\it domain 0}. The domain builder 
is responsible for building the initial page tables for a domain  
and loading its kernel image at the appropriate virtual address. 



\chapter{Network I/O}

Virtual network device services are provided by shared memory
communications with a `backend' domain.  From the point of view of
other domains, the backend may be viewed as a virtual ethernet switch
element with each domain having one or more virtual network interfaces
connected to it.

\section{Backend Packet Handling}
The backend driver is responsible primarily for {\it data-path} operations.
In terms of networking this means packet transmission and reception.

On the transmission side, the backend needs to perform two key actions:
\begin{itemize}
\item {\tt Validation:} A domain may only be allowed to emit packets
matching a certain specification; for example, ones in which the
source IP address matches one assigned to the virtual interface over
which it is sent.  The backend would be responsible for ensuring any
such requirements are met, either by checking or by stamping outgoing
packets with prescribed values for certain fields.

Validation functions can be configured using standard firewall rules
(i.e. IP Tables, in the case of Linux).

\item {\tt Scheduling:} Since a number of domains can share a single
``real'' network interface, the hypervisor must mediate access when
several domains each have packets queued for transmission.  Of course,
this general scheduling function subsumes basic shaping or
rate-limiting schemes.

\item {\tt Logging and Accounting:} The hypervisor can be configured
with classifier rules that control how packets are accounted or
logged.  For example, {\it domain0} could request that it receives a
log message or copy of the packet whenever another domain attempts to
send a TCP packet containg a SYN.
\end{itemize}

On the recive side, the backend's role is relatively straightforward:
once a packet is received, it just needs to determine the virtual interface(s)
to which it must be delivered and deliver it via page-flipping. 


\section{Data Transfer}

Each virtual interface uses two ``descriptor rings'', one for transmit,
the other for receive.  Each descriptor identifies a block of contiguous
physical memory allocated to the domain.  There are four cases:

\begin{itemize}

\item The transmit ring carries packets to transmit from the domain to the
hypervisor.

\item The return path of the transmit ring carries ``empty'' descriptors
indicating that the contents have been transmitted and the memory can be
re-used.

\item The receive ring carries empty descriptors from the domain to the 
hypervisor; these provide storage space for that domain's received packets.

\item The return path of the receive ring carries packets that have been
received.
\end{itemize}

Real physical addresses are used throughout, with the domain performing 
translation from pseudo-physical addresses if that is necessary.

If a domain does not keep its receive ring stocked with empty buffers then 
packets destined to it may be dropped.  This provides some defense against 
receiver-livelock problems because an overload domain will cease to receive
further data.  Similarly, on the transmit path, it provides the application
with feedback on the rate at which packets are able to leave the system.

Synchronization between the hypervisor and the domain is achieved using 
counters held in shared memory that is accessible to both.  Each ring has
associated producer and consumer indices indicating the area in the ring
that holds descriptors that contain data.  After receiving {\it n} packets
or {\t nanoseconds} after receiving the first packet, the hypervisor sends
an event to the domain. 

\chapter{Block I/O}

\section{Virtual Block Devices (VBDs)}

All guest OS disk access goes through the VBD interface.  The VBD
interface provides the administrator with the ability to selectively
grant domains access to portions of block storage devices visible to
the the block backend device (usually domain 0).

VBDs can literally be backed by any block device accessible to the
backend domain, including network-based block devices (iSCSI, *NBD,
etc), loopback devices and LVM / MD devices.

Old (Xen 1.2) virtual disks are not supported under Xen 2.0, since
similar functionality can be achieved using the (more advanced) LVM
system, which is already in widespread use.

\subsection{Data Transfer}
Domains which have been granted access to a logical block device are permitted
to read and write it by shared memory communications with the backend domain. 

In overview, the same style of descriptor-ring that is used for
network packets is used here.  Each domain has one ring that carries
operation requests to the hypervisor and carries the results back
again.

Rather than copying data, the backend simply maps the domain's buffers
in order to enable direct DMA to them.  The act of mapping the buffers
also increases the reference counts of the underlying pages, so that
the unprivileged domain cannot try to return them to the hypervisor,
install them as page tables, or any other unsafe behaviour.
%block API here 

\chapter{Privileged operations}
{\it Domain0} is responsible for building all other domains on the server
and providing control interfaces for managing scheduling, networking, and
blocks.

\chapter{CPU Scheduler}

Xen offers a uniform API for CPU schedulers.  It is possible to choose
from a number of schedulers at boot and it should be easy to add more.

\paragraph*{Note: SMP host support}
Xen has always supported SMP host systems.  Domains are statically assigned to
CPUs, either at creation time or when manually pinning to a particular CPU.
The current schedulers then run locally on each CPU to decide which of the
assigned domains should be run there.

\section{Standard Schedulers}

These BVT, Atropos and Round Robin schedulers are part of the normal
Xen distribution.  BVT provides proportional fair shares of the CPU to
the running domains.  Atropos can be used to reserve absolute shares
of the CPU for each domain.  Round-robin is provided as an example of
Xen's internal scheduler API.

More information on the characteristics and use of these schedulers is
available in {\tt Sched-HOWTO.txt}.




\appendix

%\newcommand{\hypercall}[1]{\vspace{5mm}{\large\sf #1}}





\newcommand{\hypercall}[1]{\vspace{2mm}{\sf #1}}






\chapter{Xen Hypercalls}
\label{a:hypercalls}

Hypercalls represent the procedural interface to Xen; this appendix 
categorizes and describes the current set of hypercalls. 

\section{Invoking Hypercalls} 

Hypercalls are invoked in a manner analogous to system calls in a
conventional operating system; a software interrupt is issued which
vectors to an entry point within Xen. On x86\_32 machines the
instruction required is {\tt int \$82}; the (real) IDT is setup so
that this may only be issued from within ring 1. The particular 
hypercall to be invoked is contained in {\tt EAX} --- a list 
mapping these values to symbolic hypercall names can be found 
in {\tt xen/include/public/xen.h}. 

On some occasions a set of hypercalls will be required to carry
out a higher-level function; a good example is when a guest 
operating wishes to context switch to a new process which 
requires updating various privileged CPU state. As an optimization
for these cases, there is a generic mechanism to issue a set of 
hypercalls as a batch: 

\begin{quote}
\hypercall{multicall(void *call\_list, int nr\_calls)}

Execute a series of hypervisor calls; {\tt nr\_calls} is the length of
the array of {\tt multicall\_entry\_t} structures pointed to be {\tt
call\_list}. Each entry contains the hypercall operation code followed
by up to 7 word-sized arguments.
\end{quote}

Note that multicalls are provided purely as an optimization; there is
no requirement to use them when first porting a guest operating
system.


\section{Virtual CPU Setup} 

At start of day, a guest operating system needs to setup the virtual
CPU it is executing on. This includes installing vectors for the
virtual IDT so that the guest OS can handle interrupts, page faults,
etc. However the very first thing a guest OS must setup is a pair 
of hypervisor callbacks: these are the entry points which Xen will
use when it wishes to notify the guest OS of an occurrence. 

\begin{quote}
\hypercall{set\_callbacks(unsigned long event\_selector, unsigned long
  event\_address, unsigned long failsafe\_selector, unsigned long
  failsafe\_address) }

Register the normal (``event'') and failsafe callbacks for 
event processing. In each case the code segment selector and 
address within that segment are provided. The selectors must
have RPL 1; in XenLinux we simply use the kernel's CS for both 
{\tt event\_selector} and {\tt failsafe\_selector}.

The value {\tt event\_address} specifies the address of the guest OSes
event handling and dispatch routine; the {\tt failsafe\_address}
specifies a separate entry point which is used only if a fault occurs
when Xen attempts to use the normal callback. 
\end{quote} 


After installing the hypervisor callbacks, the guest OS can 
install a `virtual IDT' by using the following hypercall: 

\begin{quote} 
\hypercall{set\_trap\_table(trap\_info\_t *table)} 

Install one or more entries into the per-domain 
trap handler table (essentially a software version of the IDT). 
Each entry in the array pointed to by {\tt table} includes the 
exception vector number with the corresponding segement selector 
and entry point. Most guest OSes can use the same handlers on 
Xen as when running on the real hardware; an exception is the 
page fault handler (exception vector 14) where a modified 
stack-frame layout is used. 


\end{quote} 

Finally, as an optimization it is possible for each guest OS 
to install one ``fast trap'': this is a trap gate which will 
allow direct transfer of control from ring 3 into ring 1 without
indirecting via Xen. In most cases this is suitable for use by 
the guest OS system call mechanism, although it may be used for
any purpose. 


\begin{quote}
\hypercall{set\_fast\_trap(int idx)}

Install the handler for exception vector {\tt idx} as the ``fast
trap'' for this domain. Note that this installs the current handler 
(i.e. that which has been installed more recently via a call 
to {\tt set\_trap\_table()}). 

\end{quote}



\section{Scheduling and Timer}

Domains are premptively scheduled by Xen according to the 
parameters installed by domain 0 (see Section~\ref{s:dom0ops}). 
In addition, however, a domain may choose to explicitly 
control certain behaviour with the following hypercall: 

\begin{quote} 
\hypercall{sched\_op(unsigned long op)} 

Request scheduling operation from hypervisor. The options are: {\it
yield}, {\it block}, and {\it shutdown}.  {\it yield} keeps the
calling domain runnable but may cause a reschedule if other domains
are runnable.  {\it block} removes the calling domain from the run
queue and cause is to sleeps until an event is delivered to it.  {\it
shutdown} is used to end the domain's execution; the caller can
additionally specify whether the domain should reboot, halt or
suspend.
\end{quote} 

To aid the implementation of a process scheduler within a guest OS,
Xen provides a virtual programmable timer:

\begin{quote}
\hypercall{set\_timer\_op(uint64\_t timeout)} 

Request a timer event to be sent at the specified system time (time 
in nanoseconds since system boot). The hypercall actually passes the 
64-bit timeout value as a pair of 32-bit values. 

\end{quote} 

Note that calling {\tt set\_timer\_op()} prior to {\tt sched\_op} 
allows block-with-timeout semantics. 


\section{Page Table Management} 

Since guest operating systems have read-only access to their page 
tables, Xen must be involved when making any changes. The following
multi-purpose hypercall can be used to modify page-table entries, 
update the machine-to-physical mapping table, flush the TLB, install 
a new page-table base pointer, and more.

\begin{quote} 
\hypercall{mmu\_update(mmu\_update\_t *req, int count, int *success\_count)} 

Update the page table for the domain; a set of {\tt count} updates are
submitted for processing in a batch, with {\tt success\_count} being 
updated to report the number of successfull updates.  

Each element of {\tt req[]} contains a pointer (address) and value; 
the least significant 2-bits of the pointer are used to distinguish 
the type of update requested as follows:
\begin{description} 

\item[\it MMU\_NORMAL\_PT\_UPDATE:] update a page directory entry or
page table entry to the associated value; Xen will check that the
update is safe, as described in Chapter~\ref{c:memory}.

\item[\it MMU\_MACHPHYS\_UPDATE:] update an entry in the
  machine-to-physical table. The calling domain must own the machine
  page in question (or be privileged).

\item[\it MMU\_EXTENDED\_COMMAND:] perform additional MMU operations.
The set of additional MMU operations is considerable, and includes
updating {\tt cr3} (or just re-installing it for a TLB flush),
flushing the cache, installing a new LDT, or pinning \& unpinning
page-table pages (to ensure their reference count doesn't drop to zero
which would require a revalidation of all entries).

Further extended commands are used to deal with granting and 
acquiring page ownership; see Section~\ref{s:idc}. 


\end{description}

More details on the precise format of all commands can be 
found in {\tt xen/include/public/xen.h}. 


\end{quote}

Explicitly updating batches of page table entries is extremely
efficient, but can require a number of alterations to the guest
OS. Using the writable page table mode (Chapter~\ref{c:memory}) is
recommended for new OS ports.

Regardlesss of which page table update mode is being used, however,
there are some occasions (notably handling a demand page fault) where
a guest OS will wish to modify exactly one PTE rather than a
batch. This is catered for by the following:

\begin{quote} 
\hypercall{update\_va\_mapping(unsigned long page\_nr, unsigned long
val, \\ unsigned long flags)}

Update the currently installed PTE for the page {\tt page\_nr} to 
{\tt val}. As with {\tt mmu\_update()}, Xen checks the modification 
is safe before applying it. The {\tt flags} determine which kind
of TLB flush, if any, should follow the update. 

\end{quote} 

Finally, sufficiently privileged domains may occasionally wish to manipulate 
the pages of others: 
\begin{quote}

\hypercall{update\_va\_mapping\_otherdomain(unsigned long page\_nr,
unsigned long val, unsigned long flags, uint16\_t domid)}

Identical to {\tt update\_va\_mapping()} save that the pages being
mapped must belong to the domain {\tt domid}. 

\end{quote}

This privileged operation is currently used by backend virtual device
drivers to safely map pages containing I/O data. 



\section{Segmentation Support}

Xen allows guest OSes to install a custom GDT if they require it; 
this is context switched transparently whenever a domain is 
[de]scheduled.  The following hypercall is effectively a 
`safe' version of {\tt lgdt}: 

\begin{quote}
\hypercall{set\_gdt(unsigned long *frame\_list, int entries)} 

Install a global descriptor table for a domain; {\tt frame\_list} is
an array of up to 16 machine page frames within which the GDT resides,
with {\tt entries} being the actual number of descriptor-entry
slots. All page frames must be mapped read-only within the guest's
address space, and the table must be large enough to contain Xen's
reserved entries (see {\tt xen/include/public/arch-x86\_32.h}).

\end{quote}

Many guest OSes will also wish to install LDTs; this is achieved by
using {\tt mmu\_update()} with an extended commmand, passing the
linear address of the LDT base along with the number of entries. No
special safety checks are required; Xen needs to perform this task
simply since {\tt lldt} requires CPL 0.


Xen also allows guest operating systems to update just an 
individual segment descriptor in the GDT or LDT:  

\begin{quote}
\hypercall{update\_descriptor(unsigned long ma, unsigned long word1,
unsigned long word2)}

Update the GDT/LDT entry at machine address {\tt ma}; the new
8-byte descriptor is stored in {\tt word1} and {\tt word2}.
Xen performs a number of checks to ensure the descriptor is 
valid. 

\end{quote}

Guest OSes can use the above in place of context switching entire 
LDTs (or the GDT) when the number of changing descriptors is small. 

\section{Context Switching} 

When a guest OS wishes to context switch between two processes, 
it can use the page table and segmentation hypercalls described
above to perform the the bulk of the privileged work. In addition, 
however, it will need to invoke Xen to switch the kernel (ring 1) 
stack pointer: 

\begin{quote} 
\hypercall{stack\_switch(unsigned long ss, unsigned long esp)} 

Request kernel stack switch from hypervisor; {\tt ss} is the new 
stack segement, which {\tt esp} is the new stack pointer. 

\end{quote} 

A final useful hypercall for context switching allows ``lazy'' 
save and restore of floating point state: 

\begin{quote}
\hypercall{fpu\_taskswitch(void)} 

This call instructs Xen to set the {\tt TS} bit in the {\tt cr0}
control register; this means that the next attempt to use floating
point will cause a trap which the guest OS can trap. Typically it will
then save/restore the FP state, and clear the {\tt TS} bit. 
\end{quote} 

This is provided as an optimization only; guest OSes can also choose
to save and restore FP state on all context switches for simplicity. 


\section{Physical Memory Management}

As mentioned previously, each domain has a maximum and current 
memory allocation. The maximum allocation, set at domain creation 
time, cannot be modified. However a domain can choose to reduce 
and subsequently grow its current allocation by using the
following call: 

\begin{quote} 
\hypercall{dom\_mem\_op(unsigned int op, unsigned long *extent\_list,
  unsigned long nr\_extents, unsigned int extent\_order)}

Increase or decrease current memory allocation (as determined by 
the value of {\tt op}). Each invocation provides a list of 
extents each of which is $2^s$ pages in size, 
where $s$ is the value of {\tt extent\_order}. 

\end{quote} 

In addition to simply reducing or increasing the current memory
allocation via a `balloon driver', this call is also useful for 
obtaining contiguous regions of machine memory when required (e.g. 
for certain PCI devices, or if using superpages).  


\section{Inter-Domain Communication}
\label{s:idc} 

Xen provides a simple asynchronous notification mechanism via
\emph{event channels}. Each domain has a set of end-points (or
\emph{ports}) which may be bound to an event source (e.g. a physical
IRQ, a virtual IRQ, or an port in another domain). When a pair of
end-points in two different domains are bound together, then a `send'
operation on one will cause an event to be received by the destination
domain.

The control and use of event channels involves the following hypercall: 

\begin{quote}
\hypercall{event\_channel\_op(evtchn\_op\_t *op)} 

Inter-domain event-channel management; {\tt op} is a discriminated 
union which allows the following 7 operations: 

\begin{description} 

\item[\it alloc\_unbound:] allocate a free (unbound) local
  port and prepare for connection from a specified domain. 
\item[\it bind\_virq:] bind a local port to a virtual 
IRQ; any particular VIRQ can be bound to at most one port per domain. 
\item[\it bind\_pirq:] bind a local port to a physical IRQ;
once more, a given pIRQ can be bound to at most one port per
domain. Furthermore the calling domain must be sufficiently
privileged.
\item[\it bind\_interdomain:] construct an interdomain event 
channel; in general, the target domain must have previously allocated 
an unbound port for this channel, although this can be bypassed by 
privileged domains during domain setup. 
\item[\it close:] close an interdomain event channel. 
\item[\it send:] send an event to the remote end of a 
interdomain event channel. 
\item[\it status:] determine the current status of a local port. 
\end{description} 

For more details see
{\tt xen/include/public/event\_channel.h}. 

\end{quote} 

Event channels are the fundamental communication primitive between 
Xen domains and seamlessly support SMP. However they provide little
bandwidth for communication {\sl per se}, and hence are typically 
married with a piece of shared memory to produce effective and 
high-performance inter-domain communication. 

Safe sharing of memory pages between guest OSes is carried out granting
access on a per page basis to individual domains. This is achieved 
by using the {\tt grant\_table\_op()} hypercall. 

XXX SMH: damn, found some more calls - needs little bit more yet.

\hypercall{grant\_table\_op(unsigned int cmd, void *uop, unsigned int count)}

\hypercall{physdev\_op(void *physdev\_op)}

\hypercall{vm\_assist(unsigned int cmd, unsigned int type)}


\section{Administrative Operations}
\label{s:dom0ops}

A large number of control operations are available to a sufficiently
privileged domain (typically domain 0). These allow the creation and
management of new domains, for example. A complete list is given 
below: for more details on any or all of these, please see 
{\tt xen/include/public/dom0\_ops.h} 


\begin{quote}
\hypercall{dom0\_op(dom0\_op\_t *op)} 

Administrative domain operations for domain management. The options are:

\begin{description} 
\item [\it DOM0\_CREATEDOMAIN:] create a new domain

\item [\it DOM0\_PAUSEDOMAIN:] remove a domain from the scheduler run 
queue. 

\item [\it DOM0\_UNPAUSEDOMAIN:] mark a paused domain as schedulable
  once again. 

\item [\it DOM0\_DESTROYDOMAIN:] deallocate all resources associated
with a domain

\item [\it DOM0\_GETMEMLIST:] get list of pages used by the domain

\item [\it DOM0\_SCHEDCTL:]

\item [\it DOM0\_ADJUSTDOM:] adjust scheduling priorities for domain

\item [\it DOM0\_BUILDDOMAIN:] do final guest OS setup for domain

\item [\it DOM0\_GETDOMAINFO:] get statistics about the domain

\item [\it DOM0\_GETPAGEFRAMEINFO:] 

\item [\it DOM0\_GETPAGEFRAMEINFO2:]

\item [\it DOM0\_IOPL:] set I/O privilege level

\item [\it DOM0\_MSR:] read or write model specific registers

\item [\it DOM0\_DEBUG:] interactively invoke the debugger

\item [\it DOM0\_SETTIME:] set system time

\item [\it DOM0\_READCONSOLE:] read console content from hypervisor buffer ring

\item [\it DOM0\_PINCPUDOMAIN:] pin domain to a particular CPU

\item [\it DOM0\_GETTBUFS:] get information about the size and location of
                      the trace buffers (only on trace-buffer enabled builds)

\item [\it DOM0\_PHYSINFO:] get information about the host machine

\item [\it DOM0\_PCIDEV\_ACCESS:] modify PCI device access permissions

\item [\it DOM0\_SCHED\_ID:] get the ID of the current Xen scheduler

\item [\it DOM0\_SHADOW\_CONTROL:] switch between shadow page-table modes

\item [\it DOM0\_SETDOMAININITIALMEM:] set initial memory allocation of a domain

\item [\it DOM0\_SETDOMAINMAXMEM:] set maximum memory allocation of a domain

\item [\it DOM0\_SETDOMAINVMASSIST:] set domain VM assist options
\end{description} 
\end{quote} 





\section{Debugging Hypercalls} 

A few additional hypercalls are maintly useful for debugging: 

\begin{quote} 
\hypercall{console\_io(int cmd, int count, char *str)}

Use Xen to interact with the console; operations are:

{\it CONSOLEIO\_write}: Output count characters from buffer str.

{\it CONSOLEIO\_read}: Input at most count characters into buffer str.
\end{quote} 

A pair of hypercalls allows access to the underlying debug registers: 
\begin{quote}
\hypercall{set\_debugreg(int reg, unsigned long value)}

Set debug register {\tt reg} to {\tt value} 

\hypercall{get\_debugreg(int reg)}

Return the contents of the debug register {\tt reg}
\end{quote}

And finally, a sometimes useful call is: 
\begin{quote}
\hypercall{xen\_version(int cmd)}

Request Xen version number.
\end{quote} 







%% 
%% XXX SMH: not really sure how useful below is -- if it's still 
%% actually true, might be useful for someone wanting to write a 
%% new scheduler... not clear how many of them there are...
%%

\begin{comment}

\chapter{Scheduling API}  

The scheduling API is used by both the schedulers described above and should
also be used by any new schedulers.  It provides a generic interface and also
implements much of the ``boilerplate'' code.

Schedulers conforming to this API are described by the following
structure:

\begin{verbatim}
struct scheduler
{
    char *name;             /* full name for this scheduler      */
    char *opt_name;         /* option name for this scheduler    */
    unsigned int sched_id;  /* ID for this scheduler             */

    int          (*init_scheduler) ();
    int          (*alloc_task)     (struct task_struct *);
    void         (*add_task)       (struct task_struct *);
    void         (*free_task)      (struct task_struct *);
    void         (*rem_task)       (struct task_struct *);
    void         (*wake_up)        (struct task_struct *);
    void         (*do_block)       (struct task_struct *);
    task_slice_t (*do_schedule)    (s_time_t);
    int          (*control)        (struct sched_ctl_cmd *);
    int          (*adjdom)         (struct task_struct *,
                                    struct sched_adjdom_cmd *);
    s32          (*reschedule)     (struct task_struct *);
    void         (*dump_settings)  (void);
    void         (*dump_cpu_state) (int);
    void         (*dump_runq_el)   (struct task_struct *);
};
\end{verbatim}

The only method that {\em must} be implemented is
{\tt do\_schedule()}.  However, if there is not some implementation for the
{\tt wake\_up()} method then waking tasks will not get put on the runqueue!

The fields of the above structure are described in more detail below.

\subsubsection{name}

The name field should point to a descriptive ASCII string.

\subsubsection{opt\_name}

This field is the value of the {\tt sched=} boot-time option that will select
this scheduler.

\subsubsection{sched\_id}

This is an integer that uniquely identifies this scheduler.  There should be a
macro corrsponding to this scheduler ID in {\tt <xen/sched-if.h>}.

\subsubsection{init\_scheduler}

\paragraph*{Purpose}

This is a function for performing any scheduler-specific initialisation.  For
instance, it might allocate memory for per-CPU scheduler data and initialise it
appropriately.

\paragraph*{Call environment}

This function is called after the initialisation performed by the generic
layer.  The function is called exactly once, for the scheduler that has been
selected.

\paragraph*{Return values}

This should return negative on failure --- this will cause an
immediate panic and the system will fail to boot.

\subsubsection{alloc\_task}

\paragraph*{Purpose}
Called when a {\tt task\_struct} is allocated by the generic scheduler
layer.  A particular scheduler implementation may use this method to
allocate per-task data for this task.  It may use the {\tt
sched\_priv} pointer in the {\tt task\_struct} to point to this data.

\paragraph*{Call environment}
The generic layer guarantees that the {\tt sched\_priv} field will
remain intact from the time this method is called until the task is
deallocated (so long as the scheduler implementation does not change
it explicitly!).

\paragraph*{Return values}
Negative on failure.

\subsubsection{add\_task}

\paragraph*{Purpose}

Called when a task is initially added by the generic layer.

\paragraph*{Call environment}

The fields in the {\tt task\_struct} are now filled out and available for use.
Schedulers should implement appropriate initialisation of any per-task private
information in this method.

\subsubsection{free\_task}

\paragraph*{Purpose}

Schedulers should free the space used by any associated private data
structures.

\paragraph*{Call environment}

This is called when a {\tt task\_struct} is about to be deallocated.
The generic layer will have done generic task removal operations and
(if implemented) called the scheduler's {\tt rem\_task} method before
this method is called.

\subsubsection{rem\_task}

\paragraph*{Purpose}

This is called when a task is being removed from scheduling (but is
not yet being freed).

\subsubsection{wake\_up}

\paragraph*{Purpose}

Called when a task is woken up, this method should put the task on the runqueue
(or do the scheduler-specific equivalent action).

\paragraph*{Call environment}

The task is already set to state RUNNING.

\subsubsection{do\_block}

\paragraph*{Purpose}

This function is called when a task is blocked.  This function should
not remove the task from the runqueue.

\paragraph*{Call environment}

The EVENTS\_MASTER\_ENABLE\_BIT is already set and the task state changed to
TASK\_INTERRUPTIBLE on entry to this method.  A call to the {\tt
  do\_schedule} method will be made after this method returns, in
order to select the next task to run.

\subsubsection{do\_schedule}

This method must be implemented.

\paragraph*{Purpose}

The method is called each time a new task must be chosen for scheduling on the
current CPU.  The current time as passed as the single argument (the current
task can be found using the {\tt current} macro).

This method should select the next task to run on this CPU and set it's minimum
time to run as well as returning the data described below.

This method should also take the appropriate action if the previous
task has blocked, e.g. removing it from the runqueue.

\paragraph*{Call environment}

The other fields in the {\tt task\_struct} are updated by the generic layer,
which also performs all Xen-specific tasks and performs the actual task switch
(unless the previous task has been chosen again).

This method is called with the {\tt schedule\_lock} held for the current CPU
and local interrupts disabled.

\paragraph*{Return values}

Must return a {\tt struct task\_slice} describing what task to run and how long
for (at maximum).

\subsubsection{control}

\paragraph*{Purpose}

This method is called for global scheduler control operations.  It takes a
pointer to a {\tt struct sched\_ctl\_cmd}, which it should either
source data from or populate with data, depending on the value of the
{\tt direction} field.

\paragraph*{Call environment}

The generic layer guarantees that when this method is called, the
caller selected the correct scheduler ID, hence the scheduler's
implementation does not need to sanity-check these parts of the call.

\paragraph*{Return values}

This function should return the value to be passed back to user space, hence it
should either be 0 or an appropriate errno value.

\subsubsection{sched\_adjdom}

\paragraph*{Purpose}

This method is called to adjust the scheduling parameters of a particular
domain, or to query their current values.  The function should check
the {\tt direction} field of the {\tt sched\_adjdom\_cmd} it receives in
order to determine which of these operations is being performed.

\paragraph*{Call environment}

The generic layer guarantees that the caller has specified the correct
control interface version and scheduler ID and that the supplied {\tt
task\_struct} will not be deallocated during the call (hence it is not
necessary to {\tt get\_task\_struct}).

\paragraph*{Return values}

This function should return the value to be passed back to user space, hence it
should either be 0 or an appropriate errno value.

\subsubsection{reschedule}

\paragraph*{Purpose}

This method is called to determine if a reschedule is required as a result of a
particular task.

\paragraph*{Call environment}
The generic layer will cause a reschedule if the current domain is the idle
task or it has exceeded its minimum time slice before a reschedule.  The
generic layer guarantees that the task passed is not currently running but is
on the runqueue.

\paragraph*{Return values}

Should return a mask of CPUs to cause a reschedule on.

\subsubsection{dump\_settings}

\paragraph*{Purpose}

If implemented, this should dump any private global settings for this
scheduler to the console.

\paragraph*{Call environment}

This function is called with interrupts enabled.

\subsubsection{dump\_cpu\_state}

\paragraph*{Purpose}

This method should dump any private settings for the specified CPU.

\paragraph*{Call environment}

This function is called with interrupts disabled and the {\tt schedule\_lock}
for the specified CPU held.

\subsubsection{dump\_runq\_el}

\paragraph*{Purpose}

This method should dump any private settings for the specified task.

\paragraph*{Call environment}

This function is called with interrupts disabled and the {\tt schedule\_lock}
for the task's CPU held.

\end{comment} 




%%
%% XXX SMH: we probably should have something in here on debugging 
%% etc; this is a kinda developers manual and many devs seem to 
%% like debugging support :^) 
%% Possibly sanitize below, else wait until new xendbg stuff is in 
%% (and/or kip's stuff?) and write about that instead? 
%%

\begin{comment} 

\chapter{Debugging}

Xen provides tools for debugging both Xen and guest OSes.  Currently, the
Pervasive Debugger provides a GDB stub, which provides facilities for symbolic
debugging of Xen itself and of OS kernels running on top of Xen.  The Trace
Buffer provides a lightweight means to log data about Xen's internal state and
behaviour at runtime, for later analysis.

\section{Pervasive Debugger}

Information on using the pervasive debugger is available in pdb.txt.


\section{Trace Buffer}

The trace buffer provides a means to observe Xen's operation from domain 0.
Trace events, inserted at key points in Xen's code, record data that can be
read by the {\tt xentrace} tool.  Recording these events has a low overhead
and hence the trace buffer may be useful for debugging timing-sensitive
behaviours.

\subsection{Internal API}

To use the trace buffer functionality from within Xen, you must {\tt \#include
<xen/trace.h>}, which contains definitions related to the trace buffer.  Trace
events are inserted into the buffer using the {\tt TRACE\_xD} ({\tt x} = 0, 1,
2, 3, 4 or 5) macros.  These all take an event number, plus {\tt x} additional
(32-bit) data as their arguments.  For trace buffer-enabled builds of Xen these
will insert the event ID and data into the trace buffer, along with the current
value of the CPU cycle-counter.  For builds without the trace buffer enabled,
the macros expand to no-ops and thus can be left in place without incurring
overheads.

\subsection{Trace-enabled builds}

By default, the trace buffer is enabled only in debug builds (i.e. {\tt NDEBUG}
is not defined).  It can be enabled separately by defining {\tt TRACE\_BUFFER},
either in {\tt <xen/config.h>} or on the gcc command line.

The size (in pages) of the per-CPU trace buffers can be specified using the
{\tt tbuf\_size=n } boot parameter to Xen.  If the size is set to 0, the trace
buffers will be disabled.

\subsection{Dumping trace data}

When running a trace buffer build of Xen, trace data are written continuously
into the buffer data areas, with newer data overwriting older data.  This data
can be captured using the {\tt xentrace} program in domain 0.

The {\tt xentrace} tool uses {\tt /dev/mem} in domain 0 to map the trace
buffers into its address space.  It then periodically polls all the buffers for
new data, dumping out any new records from each buffer in turn.  As a result,
for machines with multiple (logical) CPUs, the trace buffer output will not be
in overall chronological order.

The output from {\tt xentrace} can be post-processed using {\tt
xentrace\_cpusplit} (used to split trace data out into per-cpu log files) and
{\tt xentrace\_format} (used to pretty-print trace data).  For the predefined
trace points, there is an example format file in {\tt tools/xentrace/formats }.

For more information, see the manual pages for {\tt xentrace}, {\tt
xentrace\_format} and {\tt xentrace\_cpusplit}.

\end{comment} 




\end{document}
